(1) 椭圆长半轴 \(a=(r_1+r_2)/2\)，比机械能 \(\varepsilon=-\mu/(2a)\)。由 vis-viva 公式 \(v^2=\mu(2/r-1/a)\)，两端
\[
v_p=\sqrt{\mu(2/r_1-1/a)},\qquad v_a=\sqrt{\mu(2/r_2-1/a)}.
\]
初末圆轨道速率分别 \(v_{c1}=\sqrt{\mu/r_1},v_{c2}=\sqrt{\mu/r_2}\)。

(2) 首次近地点切向加速 \(\Delta v_1=v_p-v_{c1}>0\)，到远地点还需切向加速 \(\Delta v_2=v_{c2}-v_a>0\)，所以 \(\Delta v_{\rm total}=\Delta v_1+\Delta v_2\)。这不同于简单将终末两个圆轨道速率直接相减。

(3) 开普勒第三定律给椭圆周期 \(2\pi\sqrt{a^3/\mu}\)，从近日点到远日点恰半个周期，故 \(t_{\rm tr}=\pi\sqrt{a^3/\mu}\)。理想脉冲时间忽略不计，实际连续推力需要重新积分轨道。

(4) 圆轨道比机械能 \(\varepsilon_c(r)=-\mu/(2r)\)，终末与初态之差为 \(\mu(1/r_1-1/r_2)/2>0\)。两次发动机分别改变动能，沿椭圆无推力时机械能守恒；当 \(r_2\to r_1\)，两脉冲均趋零，轨道退化为原圆。
